% ---------------------------------------------------------------------
%  Quatrième de couverture : même style que la première
% ---------------------------------------------------------------------
\cleardoublepage
\thispagestyle{empty}
\null\clearpage
\thispagestyle{empty}
\begin{tikzpicture}[remember picture,overlay]
  \fill[lhblue] ([xshift=-1.1cm]current page.north east) rectangle (current page.south east);
  \node[anchor=north west,text width=14.5cm,text=lhfonce,font=\sffamily,align=justify]
    at ([xshift=2.5cm,yshift=-3.0cm]current page.north west) {%
    {\LARGE\bfseries\color{lhblue} \titreouvrage\par}\smallskip
    {\large \soustitreouvrage\par}\bigskip\bigskip
    Le \emph{Lakehouse} est devenu l'architecture de référence des plateformes de données : des
    fichiers ouverts, stockés à faible coût, auxquels une couche de métadonnées apporte les
    transactions, l'historique et le contrôle du schéma d'un entrepôt de données.\par\medskip
    Cet ouvrage propose de la comprendre en la construisant. Après une présentation de ses
    principes, quatre travaux pratiques guident le lecteur dans la réalisation d'une plateforme
    complète, sur les données réelles d'une place de marché en ligne : ingestion et qualité des
    données, modélisation dimensionnelle avec dbt, analyse de texte, de messages JSON et d'images,
    orchestration avec Airflow, exposition avec Trino et Superset, supervision et lineage, puis
    traitement en continu avec Kafka et Spark Structured Streaming.\par\medskip
    Tout fonctionne sur un ordinateur portable, avec des logiciels libres utilisés en entreprise.
    Chaque chapitre précise ses objectifs, les résultats attendus, chiffrés, et la façon de les
    interpréter.\par\medskip
    L'ouvrage s'adresse aux étudiants de master et d'école d'ingénieurs en informatique, en science
    des données et en informatique décisionnelle, ainsi qu'aux ingénieurs qui souhaitent se former
    à ces architectures.\par\bigskip\bigskip
    {\bfseries\color{lhblue} Les auteurs}\par\smallskip
    Pr.~Anouar Riad Solh et Dr.~Mohamed Kachmar.
  };
  \draw[black!25,line width=0.5pt] ([xshift=2.5cm,yshift=3.2cm]current page.south west) -- ([xshift=-3.0cm,yshift=3.2cm]current page.south east);
  \node[anchor=north west,text=black!65,font=\sffamily] at ([xshift=2.5cm,yshift=2.9cm]current page.south west) {\editeur};
  \node[anchor=north east,text=black!65,font=\sffamily\small] at ([xshift=-3.0cm,yshift=2.9cm]current page.south east) {ISBN \isbn};
\end{tikzpicture}
